\documentclass[11pt,a4paper]{article}
\usepackage{kotex}
\usepackage{amsmath,amssymb,booktabs,geometry}
\geometry{margin=25mm}
\title{I01 선검증 6 결과: 명시적 시계 후선택에서도 출처 분리가 유지되는가}
\author{RELATIONAL ATLAS}
\date{2026년 10월 4일}
\begin{document}
\maketitle

\section{실험의 위치}
선검증 3--5의 유한 샷 회로는 에너지 위상을 계 레지스터에
직접 인가했다. 이번에는 6-qubit 계와 2-qubit 시계를 실제
Qiskit 회로의 별도 레지스터로 준비했다. 사전 고정한 조건과
검정은 선검증 6 프로토콜, 첫 결과를 본 뒤에 고정한 반복
기준은 선검증 6b 프로토콜을 따른다. 이 실험의 질문은
\emph{같은 현상이 명시적 시계의 후선택 및 그 자원 비용 아래에서도
관측되는가}이지, 차원 3의 자발적 선택을 입증하는가가 아니다.

\section{회로와 정확 확률}
계의 운동량 $k_i\in\{-1,0,1\}$을 축당 2 qubit으로 부호화하고
$E_0(k)=\sum_i k_i^2$를 두 시계 qubit에 기록하여
\begin{equation}
  |\Psi\rangle=\sum_k\sqrt{w_k}\,|k\rangle|E_0(k)\rangle
\end{equation}
을 준비했다. 시계의 대각 위상은
$\exp[i\tau(E_0+0.37E_0^2/3)]$, 계의 상대 위상은
$\exp(i k\cdot\theta)$이다. 시계에 $H^{\otimes2}$를 적용하고
$00$을 수락하며, 계는 폭 $\beta_r=0.15$인 기준 상태의
역준비 뒤 $0^6$을 성공으로 센다. 모든 각도에서
\begin{equation}
  P(C=00)=\tfrac14,\qquad
  P(C=00,S=0^6)=\tfrac14 p_{\rm direct}(\theta,\tau).
\end{equation}
이는 물리적인 시계가 스스로 만들어진다는 주장이 아니라
\emph{에너지 상관 시계 상태를 준비할 수 있다}는 가정 아래의
회로 동치다. 네 대표 조건에 대한 Statevector/Aer density-matrix의
최대 절대 오차는 $1.41\times10^{-14}$였고, 합성 회로 깊이는
588--609였다. 따라서 회로 구현이 정확 모형의 동시확률을
재현한다는 검사를 통과했다.

동시 성공을 1회 준비에 대한 Bernoulli 결과로 분석했다.
후선택된 약 $1/4$만 분모에 넣어 준비 비용을 숨기지 않았다.
정확 Fisher 행렬의 rank는 비제약 상태에서 표지 보존 3,
균일 표지 삭제 2, 비균일 표지 삭제 3이다.
$P=0$ 상태는 이 세 경우 모두 2다. 시계 수락이 확률을
$1/4$로 줄여 정보량은 약화되지만 이 모형의 정확 영공간은
바꾸지 않는다.

\section{유한 샷 결과와 독립 반복}
$n=3$, 표지 7개, $\tau=0,\pi/2$, 공통 과학 위상
$0\to0.6$ rad, 비제약/$P=0$, 균일/비균일 표지를 교차했다.
각 실행은 336회로, 44,040,192회 준비였으며 2회 실행의
합계는 88,080,384회 준비다. 첫 실행의 시계 수락률은
0.2499663, 두 실행 합산 수락률은 0.2499786으로 예측
0.25와 일치한다. 독립 seed는 20261027/20261029다.

사전 지정된 약한 pooled-null 검정 8개의 bootstrap $p$값은
다음과 같다. 가족별 Bonferroni 문턱은 0.00625다.
\begin{center}
\begin{tabular}{lllrrr}
\toprule
표지 분포 & 상태 & $\tau/\pi$ & 첫 실행 & 독립 반복 & 합산\\
\midrule
균일 & 비제약 & 0   & 0.0116 & 0.7998 & 0.3263\\
균일 & 비제약 & 1/2 & 0.0362 & 0.2479 & 0.8346\\
균일 & $P=0$ & 0   & 0.8263 & 0.0366 & 0.0974\\
균일 & $P=0$ & 1/2 & 0.9089 & 0.2054 & 0.3910\\
비균일 & 비제약 & 0   & 0.0053 & 0.0078 & 0.0001\\
비균일 & 비제약 & 1/2 & 0.2013 & 0.0058 & 0.0039\\
비균일 & $P=0$ & 0   & 0.8029 & 0.6826 & 0.6795\\
비균일 & $P=0$ & 1/2 & 0.6799 & 0.7512 & 0.7422\\
\bottomrule
\end{tabular}
\end{center}
첫 실행에서 미리 예상한 $\tau=0$ 양성 대조만 통과했다.
반복에서는 $\tau=0$이 문턱 바로 위였고, 약한
$\tau=\pi/2$가 문턱 아래였다. 사전 지정한 합산 주 판정은
$\tau=0$을 검출하고 정확 null 여섯 조건을 검출하지 않아
통과했다. 합산에서 약한 $\tau=\pi/2$도 문턱 아래였지만,
단독 실행끼리 검출 여부가 바뀐 만큼 \emph{견고한 독립 재현}
이라고 부르지 않는다. 첫 실행의 균일 표지 명목 $p<0.05$
두 개는 독립 반복 및 합산에서 사라졌다.

정확 모형의 동시 성공 pooled 차이는 비균일ㆍ비제약에서
$\tau=0$일 때 축당 $+0.000725$, $\tau=\pi/2$일 때
$-0.000124$다. 합산 관측 차이는 축별로 각각
$(0.000579,0.000492,0.000425)$,
$(-0.000234,-0.000246,-0.000180)$이었다.
첫 번째는 예측보다 작고 두 번째는 예측보다 큰 쪽으로
흔들렸으며, 현재 샷의 통계적 오차가 결과 해석에 중요하다.

\section{사후 추론 품질 감사}
회로별 동시 성공의 Pearson 표준화 잔차 제곱합은 첫/둘째
실행에서 333.8/350.5 (각각 336회로), 시계 수락은
298.9/330.3이었다. 이는 이항 크기의 변동과 양립하는
진단값이다. 공통 job seed 및 동시 성공/시계 수락의 상관을
무시하므로 단일 $\chi^2$ 유의성 검정으로 사용하지 않는다.

정확 동시확률에서 별도로 생성한 이항 데이터로 200패널씩
검정 절차를 \emph{사후} 감사했다. 여섯 정확 null 중 하나라도
가족 문턱을 넘은 패널은 단일/합산 예산에서 2/200, 6/200이었다.
강한 $\tau=0$ 양성의 검출은 194/200, 200/200이었다.
약한 $\tau=\pi/2$ 양성은 11/200, 15/200뿐이었다.
따라서 합산의 약한 양성 검출은 정확 모형 아래서도
가능하지만 흔치 않은 상방 변동이며, 이를 안정적 검정력의
증거로 확대하지 않는다. 200패널의 표본 오차 및 999회
bootstrap 분해능도 남아 있다.

\section{판정과 다음 경계}
명시적 시계 레지스터에서도 \emph{비균일 표지 삭제가 비제약
상태의 공통 방향을 남기지만, $P=0$ 또는 균일 표지는 이를
남기지 않는다}는 정확 구조와 강한 조건의 유한 샷 증거가
유지됐다. 후선택 수락률 $1/4$은 준비 예산의 실질적 비용이다.
하지만 이번 시계는 준비된 에너지 상관 ancilla이고 검출도
이상적인 것이다. 자발적 차원 선택, 자율 시계, 미지 표지
분포 또는 오차 채널에 대한 결론은 아니다. 발전연구의
주제는 계속 보류한다. 다음 검증은 이 회로에 구체적인
준비/위상/측정 오류를 넣어 구조적인 영과 약한 양수를
얼마나 쉽게 혼동하는지 분리하는 방향이 적절하다.
\end{document}
